% ============================================================================
% PLANTILLA 06 — LABORATORIO / PRÁCTICA CON SOFTWARE
% ----------------------------------------------------------------------------
% Para: laboratorios de cómputo, uso de software (R, Stata, Python, EViews,
% SPSS), simulaciones y programación. Integra código, salidas de consola y
% preguntas de interpretación.
% Compilar: latexmk -pdf plantilla-06-laboratorio.tex
% Solucionario: \documentclass[soluciones]{evaluacion}
% ============================================================================
\documentclass{evaluacion}

\universidad{Universidad Nacional de San Cristóbal de Huamanga}
\facultad{Facultad de Ciencias Económicas, Administrativas y Contables}
\escuela{Escuela Profesional de Economía}
\curso{Econometría II}
\codigocurso{EC-453}
\docente{Mg. Nombre Apellido}
\periodo{Semestre 2026-I}
\tipoevaluacion{Laboratorio 04}
\fechaevaluacion{21 de julio de 2026}
\duracion{2 horas}
\modalidad{Laboratorio · individual}

\begin{document}
\membrete

\begin{instrucciones}
\begin{itemize}
  \item Trabaje en la carpeta \texttt{lab04/} entregada por el docente;
        la base de datos es \texttt{salarios.dta} ($n = 526$).
  \item Entregue al final: el \emph{do-file} (o script) comentado y un
        documento con las salidas y sus interpretaciones.
  \item Se califica tanto la correcta ejecución del código como la
        calidad de la interpretación económica.
\end{itemize}
\end{instrucciones}

\begin{datos}[Descripción de variables]
\begin{tabularx}{\linewidth}{@{}l X@{}}
  \texttt{wage}  & salario por hora, en dólares\\
  \texttt{educ}  & años de educación\\
  \texttt{exper} & años de experiencia laboral\\
  \texttt{tenure}& años en el empleo actual\\
  \texttt{female}& 1 si es mujer, 0 en otro caso\\
\end{tabularx}
\end{datos}

% ============================================================ ACTIVIDADES ==
\begin{ejercicio}[4][Preparación de datos]
Cargue la base, genere la variable \texttt{lwage} (logaritmo natural del
salario) y reporte estadísticos descriptivos de todas las variables.
Escriba el código empleado.

\recuadrorespuesta{4.5cm}
\begin{solucion}
En Stata:
\begin{codigo}[language=Stata]
use "salarios.dta", clear
generate lwage = ln(wage)
summarize wage lwage educ exper tenure female
\end{codigo}
Equivalente en R:
\begin{codigo}[language=R]
datos <- haven::read_dta("salarios.dta")
datos$lwage <- log(datos$wage)
summary(datos[, c("wage", "lwage", "educ", "exper", "tenure", "female")])
\end{codigo}
\end{solucion}
\end{ejercicio}

\begin{ejercicio}[6][Estimación e interpretación]
Estime por MCO la ecuación de Mincer
\[
  \texttt{lwage} = \beta_0 + \beta_1\,\texttt{educ}
  + \beta_2\,\texttt{exper} + \beta_3\,\texttt{tenure} + u
\]
y responda a partir de la salida siguiente:

\begin{salida}
      Source |       SS           df       MS      Number of obs   =       526
-------------+----------------------------------   F(3, 522)       =     80.39
       Model |  46.8741776         3  15.6247259   Prob > F        =    0.0000
    Residual |  101.455574       522  .194359337   R-squared       =    0.3160
-------------+----------------------------------   Adj R-squared   =    0.3121
       Total |  148.329751       525  .282532860   Root MSE        =    .44086

       lwage | Coefficient  Std. err.      t    P>|t|
-------------+--------------------------------------------
        educ |   .0920290   .0073299    12.56   0.000
       exper |   .0041211   .0017233     2.39   0.017
      tenure |   .0220672   .0030936     7.13   0.000
       _cons |   .2843595   .1041904     2.73   0.007
\end{salida}

\begin{apartados}
  \item Interprete el coeficiente de \texttt{educ} en términos
        porcentuales. \pts{2}
  \item Evalúe la significancia individual de cada regresor al 5\,\%. \pts{2}
  \item Interprete el $R^{2}$ ajustado. \pts{2}
\end{apartados}
\lineasrespuesta{8}
\begin{solucion}
\textbf{a)} Un año adicional de educación se asocia con un salario por
hora mayor en aproximadamente $9{,}2\,\%$
($100\cdot\hat\beta_1$, aproximación log-nivel), manteniendo constantes
experiencia y antigüedad.

\textbf{b)} Los tres regresores son individualmente significativos al
5\,\% ($p = 0{,}000$; $0{,}017$; $0{,}000 < 0{,}05$); la constante también
lo es ($p=0{,}007$).

\textbf{c)} El modelo explica el $31{,}2\,\%$ de la variabilidad del
log-salario una vez penalizado el número de regresores — habitual en
microdatos de corte transversal, donde la heterogeneidad individual no
observada es amplia.
\end{solucion}
\end{ejercicio}

\begin{ejercicio}[6][Diagnóstico]
Contraste la presencia de heterocedasticidad con la prueba de
Breusch--Pagan y, de ser necesario, reestime con errores estándar
robustos. Escriba los comandos, reporte el valor-$p$ y concluya.

\recuadrorespuesta{5cm}
\begin{solucion}
\begin{codigo}[language=Stata]
regress lwage educ exper tenure
estat hettest            // Breusch-Pagan: chi2(1) = 15.15, p = 0.0001
regress lwage educ exper tenure, vce(robust)
\end{codigo}
Con $p = 0{,}0001 < 0{,}05$ se rechaza homocedasticidad; los errores
estándar clásicos son inválidos. Al reestimar con \texttt{vce(robust)}
los coeficientes no cambian (MCO sigue siendo insesgado) pero la
inferencia se corrige; \texttt{educ} y \texttt{tenure} permanecen
significativos al 1\,\%.
\end{solucion}
\end{ejercicio}

\begin{ejercicio}[4][Simulación]
Escriba un programa que simule 1\,000 muestras de tamaño 50 del modelo
$y = 1 + 0{,}5x + u$, con $u\sim N(0,1)$, guarde los $\hat\beta_1$ y
verifique empíricamente el insesgamiento del estimador MCO.

\recuadrorespuesta{5cm}
\begin{solucion}
\begin{codigo}[language=R]
set.seed(123)
b1 <- replicate(1000, {
  x <- rnorm(50, mean = 5, sd = 2)
  y <- 1 + 0.5 * x + rnorm(50)
  coef(lm(y ~ x))[2]
})
mean(b1)   # 0.5004: promedio ~ 0.5, consistente con insesgamiento
sd(b1)     # dispersion muestral del estimador
hist(b1, main = "Distribución de beta_1 estimado")
\end{codigo}
El promedio de los mil estimados ($\approx 0{,}500$) coincide con el
parámetro poblacional: evidencia de que $E(\hat\beta_1)=\beta_1$.
\end{solucion}
\end{ejercicio}

\begin{criterios}
\begin{tabularx}{\linewidth}{@{}X c@{}}
  \textbf{Criterio} & \textbf{Puntaje}\\
  \arrayrulecolor{grislinea}\hline\addlinespace[3pt]
  Código correcto, reproducible y comentado        & 8\\
  Lectura correcta de salidas computacionales      & 6\\
  Interpretación económica de los resultados       & 4\\
  Orden y presentación del entregable              & 2\\
\end{tabularx}
\end{criterios}

\findelexamen
\end{document}
